\section{Why no equivalent norm is AUC}\label{sec:no-auc}

The renorming obstruction uses bounded root paths and weakly null
child increments. We first isolate the elementary finite-codimensional
argument; its midpoint counterpart will be proved separately.

\begin{lemma}[Modulus conventions]\label{lem:modulus-conventions}
For an infinite-dimensional real Banach space $Y$ with norm $N$, let
$m_N$ be the modulus obtained from \eqref{eq:average-modulus} by replacing
the average by $\max\{N(x+ty),N(x-ty)\}$. Then, for $t>0$,
\[
 0\le\overline\delta_N(t)\le\widehat\delta_N(t)
 \le m_N(t)\le2\widehat\delta_N(t).
\]
For either symmetric modulus the inner infimum is unchanged if
$N(y)=1$ is replaced by $N(y)\ge1$. Both symmetric moduli, including
their centerwise versions before taking the infimum over $x$, are
nondecreasing. Positivity on $(0,1)$ therefore implies positivity at
every positive radius.
\end{lemma}

\begin{proof}
For fixed $x,y$, the average and maximum of the two signed norms are
even convex functions of the radius, hence nondecreasing on $[0,\infty)$.
This proves the assertions about radii and monotonicity. A norming
functional for a unit center $x$ makes both signed norms at least one
on its kernel. Intersecting any chosen finite-codimensional subspace
with this kernel proves nonnegativity of the one-sided modulus and
that the maximum excess is at most twice the average excess, after
taking the supremum over subspaces. The average is pointwise at most
the maximum. Finally, on a fixed subspace the infimum of the one-sided
excess is a lower bound for each of the two signed excesses, since
$y$ and $-y$ both belong to that subspace. Taking the successive
suprema and infima proves the remaining comparison.
\end{proof}

\begin{lemma}[Bounded trees]\label{lem:bounded-tree}
Let $Y$ be a real Banach space. Suppose that for every integer $h\ge1$
there are vectors $(u_s^{(h)})_{s\in\mathbb N^{\le h}}$ and constants
$m,M,c>0$, independent of $h$, such that
\[
 m\le\|u_s^{(h)}\|\le M,
 \qquad
 \|u_{s^\frown j}^{(h)}-u_s^{(h)}\|\ge c\quad(|s|<h).
\]
For every nonterminal $s$, suppose the child increments tend weakly
to zero as $j\to\infty$. Then every equivalent norm satisfying
$\alpha\|\cdot\|\le N\le\beta\|\cdot\|$, with
$0<\alpha\le\beta<\infty$, obeys
\begin{equation}\label{eq:bounded-tree-scale}
 \overline\delta_N\!\left(\frac{\alpha c}{2\beta M}\right)=0.
\end{equation}
The same conclusion holds for one infinite tree with these properties.
\end{lemma}

\begin{proof}
Put $t_0=\alpha c/(2\beta M)$. If the displayed modulus is positive,
choose $0<\gamma<\overline\delta_N(t_0)$. At a current vector
$u=u_s^{(h)}$, the definition supplies a finite-codimensional $F$ such that
\[
 N(u+t_0N(u)y)\ge(1+\gamma)N(u)
 \qquad(y\in F,\ N(y)=1).
\]
Convexity extends this to every $w\in F$ with $N(w)\ge t_0N(u)$:
for a unit center $z$ and $r\ge t_0$,
\[
 N(z+ry)-1\ge\frac r{t_0}\bigl(N(z+t_0y)-1\bigr).
\]
The weakly null increments $d_j=u_{s^\frown j}^{(h)}-u$ have images
tending to zero in norm in the finite-dimensional quotient $Y/F$.
Thus there are $w_j\in F$ with $N(w_j-d_j)\to0$. Eventually
$N(w_j)\ge\alpha c/2\ge t_0N(u)$. Taking the approximation error
at most $\gamma N(u)/2$ gives a child with
\[
 N(u_{s^\frown j}^{(h)})\ge(1+\gamma/2)N(u_s^{(h)}).
\]
The gain is uniform although $F$ may depend on the current node.
Choose $h$ with $\alpha m(1+\gamma/2)^h>\beta M$ and make this
choice successively along a branch of that tree. The resulting norm
exceeds $\beta M$, a contradiction. The increments in particular
force $Y$ to be infinite dimensional, and Lemma~\ref{lem:modulus-conventions}
gives the reverse nonnegative inequality in \eqref{eq:bounded-tree-scale}.
An infinite tree may be restricted to any sufficiently large height.
\end{proof}

We apply the lemma to the forest. For a node $v$ in the component with
root $r$, put $p_v=\sum_{r\preceq w\preceq v}e_w^*$. The root path
is one segment test, so $\|p_v\|\le1$; evaluation at a coordinate
on that path gives equality. At a nonterminal node the child increments
are $e_{v^\frown j}^*$, each of norm one. They tend weakly to zero:
$X^*=J$, and every $u\in J$ has square-summable coordinates, so
$u(v^\frown j)\to0$. We therefore have $m=M=c=1$ in
Lemma~\ref{lem:bounded-tree}. For every equivalent norm
$\alpha\|\cdot\|\le N\le\beta\|\cdot\|$,
\[
 \overline\delta_N(\alpha/(2\beta))=0.
\]
This proves the renorming assertion in Theorem~\ref{thm:main}.
